\section{Appendix refinement: isolated secondary spectral edge}
\label{sec:secondary-term}

This appendix proves a conditional finite-dimensional refinement of the
error term in \Cref{thm:class-mertens-explicit}. Main Theorem~I supplies
the Adams--M\"obius primitive-channel recovery, and Main Theorem~II
supplies the class Mertens and Adams-corrected Euler-product constants.
Under the explicit hypotheses \(\lambda>1\), \(\eta<\lambda\), and
\(R>\lambda^{1/2}\), a non-Perron spectral edge is isolated above the
square-root contribution from repeated primitive orbits, and the same
Peter--Weyl trace formula identifies the first oscillating remainder.
Without this separation hypothesis, this section asserts no secondary
expansion. The connection with the main theorem chain is the already-proved
Peter--Weyl primitive trace formula: the edge skeleton below records
exactly what that formula sees when a secondary spectral edge is isolated.

For every finite matrix \(M\) and every eigenvalue
\(\alpha\in\spec(M)\), let \(m_M(\alpha)\) denote the algebraic
multiplicity of \(\alpha\). Assume throughout this section that
\(\lambda>1\), that
\[
  \eta:=\max_{\varrho\neq\mathbf 1}\rho(A_\varrho)<\lambda,
  \qquad
  R:=\max\{\Lambda(A),\eta\}>\lambda^{1/2}.
\]
If \(G\) is trivial, the maximum defining \(\eta\) is empty and
\(\eta=0\).
Define the spectral edge sets
\[
  \begin{aligned}
    \mathcal E_{\mathbf 1}
    &:=\{\mu\in\spec(A)\setminus\{\lambda\}:|\mu|=R\}, \\
    \mathcal E_\varrho
    &:=\{\beta\in\spec(A_\varrho):|\beta|=R\}
    \quad (\varrho\neq\mathbf 1),
  \end{aligned}
\]
and, with empty maxima interpreted as \(0\),
\[
  \Gamma_{\mathrm{sec}}
  :=
  \max\left\{
    \lambda^{1/2},
    \max_{\mu\in\spec(A)\setminus(\{\lambda\}\cup\mathcal E_{\mathbf 1})}
      |\mu|,
    \max_{\varrho\neq\mathbf 1}
      \max_{\beta\in\spec(A_\varrho)\setminus\mathcal E_\varrho}
        |\beta|
  \right\}.
\]
For each conjugacy class \(C\subset G\), set
\begin{equation}\label{eq:class-secondary-profile}
  \mathfrak s_C(n)
  :=
  \frac{|C|}{|G|}
  \left(
    \sum_{\mu\in\mathcal E_{\mathbf 1}}m_A(\mu)\mu^n
    +
    \sum_{\varrho\neq\mathbf 1}\chi_\varrho^\ast(C)
      \sum_{\beta\in\mathcal E_\varrho}m_{A_\varrho}(\beta)\beta^n
  \right)
\end{equation}
and
\begin{equation}\label{eq:class-secondary-tail}
  \mathfrak T_C(N)
  :=
  \sum_{n>N}\frac{\mathfrak s_C(n)}{n\lambda^n}.
\end{equation}

\begin{theorem}[Finite edge skeleton under
  secondary separation]
  \label{thm:secondary-edge-skeleton}
  Under the assumptions of this section, define
  \[
    E_{\mathbf 1}(n)
    :=\sum_{\mu\in\mathcal E_{\mathbf 1}}m_A(\mu)\mu^n,
    \qquad
    E_\varrho(n)
    :=\sum_{\beta\in\mathcal E_\varrho}
      m_{A_\varrho}(\beta)\beta^n
    \quad(\varrho\neq\mathbf 1),
  \]
  and let
  \[
    \mathfrak b_n(C):=\frac{|G|}{|C|}\mathfrak s_C(n)
  \]
  be the normalised edge profile on conjugacy classes. Then, for every
  \(n\ge1\), the class profile is exactly the Peter--Weyl transform of
  the finite edge skeleton:
  \begin{equation}\label{eq:secondary-edge-forward}
    \mathfrak b_n(C)
    =
    E_{\mathbf 1}(n)+
    \sum_{\varrho\neq\mathbf 1}\chi_\varrho^\ast(C)E_\varrho(n).
  \end{equation}
  Conversely the edge skeleton is recovered from the normalised class
  profile by
  \begin{equation}\label{eq:secondary-edge-inverse}
    E_{\mathbf 1}(n)
    =\sum_{C\subset G}\frac{|C|}{|G|}\mathfrak b_n(C),
    \qquad
    E_\pi(n)
    =\sum_{C\subset G}\frac{|C|}{|G|}\mathfrak b_n(C)\chi_\pi(C)
    \quad(\pi\neq\mathbf 1).
  \end{equation}
  Hence this appendix secondary correction is completely exhausted by
  the finite list of edge exponential sums \(E_\varrho(n)\). In
  particular, if \(\mathfrak s_C(n)=0\) for every conjugacy class \(C\)
  and every \(n\ge1\), then all edge sums \(E_\varrho(n)\) vanish
  identically; equivalently, every set \(\mathcal E_\varrho\) is empty.
\end{theorem}

\begin{proof}
  Equation \eqref{eq:secondary-edge-forward} is the definition of
  \(\mathfrak s_C(n)\), after multiplying by \(|G|/|C|\) and writing the
  two displayed edge sums as \(E_{\mathbf1}(n)\) and \(E_\varrho(n)\).
  For the inverse formulas, average
  \eqref{eq:secondary-edge-forward} over conjugacy classes with weights
  \(|C|/|G|\). The non-trivial character averages vanish, so the first
  formula in \eqref{eq:secondary-edge-inverse} follows. Multiplying
  \eqref{eq:secondary-edge-forward} by \(\chi_\pi(C)\), averaging over
  conjugacy classes, and using the Hermitian character orthogonality
  relation
  \[
    \sum_{C\subset G}\frac{|C|}{|G|}
      \chi_\varrho^\ast(C)\chi_\pi(C)=\delta_{\varrho\pi}
  \]
  gives the second formula.

  The preceding identities show that the class data and the edge
  skeleton determine one another for each fixed \(n\). If
  \(\mathfrak s_C(n)=0\) for every \(C\) and \(n\), then
  \eqref{eq:secondary-edge-inverse} gives \(E_\varrho(n)=0\) for every
  \(\varrho\) and every \(n\). For a fixed \(\varrho\), this is a finite
  exponential sum
  \(\sum_{\alpha\in\mathcal E_\varrho}m_{A_\varrho}(\alpha)\alpha^n\)
  that vanishes for all positive \(n\). The Vandermonde matrix attached
  to the distinct non-zero frequencies in \(\mathcal E_\varrho\) is
  invertible, so all coefficients must be zero. The coefficients are
  algebraic multiplicities; therefore \(\mathcal E_\varrho\) is empty.
\end{proof}

\begin{theorem}[Secondary term above the
  square-root barrier]
  \label{thm:class-mertens-secondary-term}
  Under the assumptions above, \(\Gamma_{\mathrm{sec}}<R\), the series
  \(\mathfrak T_C(N)\) converges absolutely for every \(C\), and for
  every conjugacy class \(C\subset G\),
  \begin{equation}\label{eq:class-secondary-pointwise}
    p_{n,C}
    =
    \frac{|C|}{|G|}\frac{\lambda^n}{n}
    +\frac{\mathfrak s_C(n)}{n}
    +O\!\left(\frac{\Gamma_{\mathrm{sec}}^n}{n}\right).
  \end{equation}
  Moreover,
  \begin{equation}\label{eq:class-secondary-harmonic}
    \sum_{n\le N}\frac{p_{n,C}}{\lambda^n}
    =
    \frac{|C|}{|G|}H_N
    +\mathcal M_C(A,\tau;G)-\frac{|C|}{|G|}\gamma
    -\mathfrak T_C(N)
    +O\!\left(\frac{(\Gamma_{\mathrm{sec}}/\lambda)^N}{N}\right).
  \end{equation}
  The implied constants are uniform in the conjugacy class \(C\), in
  \(n\), and in \(N\).
\end{theorem}

\begin{proof}
  Since \(R>\lambda^{1/2}\) and the eigenvalue sets are finite, the
  definition of \(R\) implies \(\Gamma_{\mathrm{sec}}<R\).

  The trace of a finite matrix equals the sum of its eigenvalues with
  algebraic multiplicity. Therefore
  \[
    \Tr(A^n)
    =
    \lambda^n
    +\sum_{\mu\in\mathcal E_{\mathbf 1}}m_A(\mu)\mu^n
    +O(\Gamma_{\mathrm{sec}}^n),
  \]
  and, for every non-trivial \(\varrho\in\Irr(G)\),
  \[
    \Tr(A_\varrho^n)
    =
    \sum_{\beta\in\mathcal E_\varrho}m_{A_\varrho}(\beta)\beta^n
    +O(\Gamma_{\mathrm{sec}}^n).
  \]
  The implied constants are uniform in \(\varrho\) because \(\Irr(G)\)
  is finite.

  Inserting the first identity into \eqref{eq:primitive-mobius} yields
  \[
    p_n(A)
    =
    \frac{\lambda^n}{n}
    +\frac{1}{n}\sum_{\mu\in\mathcal E_{\mathbf 1}}m_A(\mu)\mu^n
    +O\!\left(\frac{\Gamma_{\mathrm{sec}}^n}{n}\right),
  \]
  because the divisor \(d=1\) contributes the spectral-edge sum displayed
  above. The divisors \(d>1\) contribute
  \[
    O\!\left(
      \frac{1}{n}\sum_{\substack{d\mid n\\ d>1}}\lambda^{n/d}
    \right)
    =
    O\!\left(\frac{\lambda^{n/2}}{n}\right)
    =
    O\!\left(\frac{\Gamma_{\mathrm{sec}}^n}{n}\right).
  \]

  Fix next \(\varrho\neq\mathbf 1\). By
  \Cref{cor:primitive-peter-weyl-trace},
  \[
    \pi_{n,\varrho}(A,\tau)
    =
    \frac{1}{n}\Tr(A_\varrho^n)
    +\frac{1}{n}
      \sum_{\substack{m\mid n\\ m>1}}\mu(m)
        \sum_{\kappa\in\Irr(G)}
          a_{\varrho,\kappa}^{(m)}\Tr(A_\kappa^{n/m}).
  \]
  The first term gives the contribution from \(\mathcal E_\varrho\). For
  the divisor sum, the same argument as in
  \Cref{lem:nontrivial-peter-weyl-channel-bound}, using
  \Cref{lem:adams-coefficients-bounded} and
  \Cref{prop:twisted-spectral-radius-bound}, yields
  \[
    \sum_{\substack{m\mid n\\ m>1}}
      \left|
        \sum_{\kappa\in\Irr(G)}
          a_{\varrho,\kappa}^{(m)}\Tr(A_\kappa^{n/m})
      \right|
    =
    O(\lambda^{n/2}),
  \]
  uniformly in \(\varrho\). Hence
  \[
    \pi_{n,\varrho}(A,\tau)
    =
    \frac{1}{n}
      \sum_{\beta\in\mathcal E_\varrho}m_{A_\varrho}(\beta)\beta^n
    +O\!\left(\frac{\Gamma_{\mathrm{sec}}^n}{n}\right)
  \]
  uniformly for all non-trivial \(\varrho\).

  Substituting the last two expansions into
  \Cref{lem:class-primitive-trivial-split} gives
  \eqref{eq:class-secondary-pointwise}. The uniformity in \(C\) follows
  from the finiteness of the conjugacy classes and the bounds
  \(|\chi_\varrho(C)|\le \dim\varrho\).

  Dividing \eqref{eq:class-secondary-pointwise} by \(\lambda^n\) yields
  \[
    \frac{p_{n,C}}{\lambda^n}
    =
    \frac{|C|}{|G|}\frac{1}{n}
    +\frac{\mathfrak s_C(n)}{n\lambda^n}
    +O\!\left(\frac{(\Gamma_{\mathrm{sec}}/\lambda)^n}{n}\right).
  \]
  Since \(R<\lambda\), the definition of \(\mathfrak s_C(n)\) shows
  that \(\mathfrak T_C(N)\) converges absolutely. Summing the preceding
  identity over \(n>N\) and applying
  \Cref{lem:geometric-harmonic-tail} gives
  \[
    \sum_{n>N}
      \left(
        \frac{p_{n,C}}{\lambda^n}
        -\frac{|C|}{|G|}\frac{1}{n}
      \right)
    =
    \mathfrak T_C(N)
    +O\!\left(\frac{(\Gamma_{\mathrm{sec}}/\lambda)^N}{N}\right).
  \]
  On the other hand, \Cref{thm:class-mertens-explicit} implies
  \[
    \mathcal M_C(A,\tau;G)-\frac{|C|}{|G|}\gamma
    =
    \sum_{n\ge 1}
      \left(
        \frac{p_{n,C}}{\lambda^n}
        -\frac{|C|}{|G|}\frac{1}{n}
      \right).
  \]
  Subtracting the tail beyond \(N\) from this identity yields
  \eqref{eq:class-secondary-harmonic}.
\end{proof}

\begin{corollary}[Detection of the dominant
  secondary modulus]
  \label{cor:class-secondary-term-detects-spectrum}
  Under the hypotheses of \Cref{thm:class-mertens-secondary-term},
  \[
    R
    =
    \max_{C\subset G}
      \limsup_{n\to\infty}
      \left|
        n\,p_{n,C}-\frac{|C|}{|G|}\lambda^n
      \right|^{1/n},
  \]
  where the maximum is taken over the conjugacy classes of \(G\).
\end{corollary}

\begin{proof}
  By \eqref{eq:class-secondary-pointwise},
  \[
    n\,p_{n,C}-\frac{|C|}{|G|}\lambda^n
    =
    \mathfrak s_C(n)+O(\Gamma_{\mathrm{sec}}^n)
    \qquad (n\to\infty),
  \]
  so
  \[
    \limsup_{n\to\infty}
      \left|
        n\,p_{n,C}-\frac{|C|}{|G|}\lambda^n
      \right|^{1/n}
    \le R
  \]
  for every \(C\).

  Assume that the displayed maximum were strictly smaller than \(R\).
  Since the set of conjugacy classes is finite, there would then exist
  \(\Gamma_0<R\) such that
  \[
    n\,p_{n,C}-\frac{|C|}{|G|}\lambda^n
    =O(\Gamma_0^n)
    \qquad\text{uniformly in } C.
  \]
  Together with \eqref{eq:class-secondary-pointwise}, this would imply
  \[
    \mathfrak s_C(n)=O(\Gamma_1^n)
    \qquad\text{for every } C,
  \]
  where \(\Gamma_1:=\max\{\Gamma_0,\Gamma_{\mathrm{sec}}\}<R\).
  Multiplying by \(|G|/|C|\) gives the class function
  \[
    \mathfrak b_n(C)
    :=
    \frac{|G|}{|C|}\mathfrak s_C(n)
    =
    \sum_{\mu\in\mathcal E_{\mathbf 1}}m_A(\mu)\mu^n
    +
    \sum_{\varrho\neq\mathbf 1}\chi_\varrho^\ast(C)
      \sum_{\beta\in\mathcal E_\varrho}m_{A_\varrho}(\beta)\beta^n,
  \]
  which would also satisfy \(\mathfrak b_n(C)=O(\Gamma_1^n)\)
  uniformly in \(C\).

  Averaging over conjugacy classes yields
  \[
    \sum_{C\subset G}\frac{|C|}{|G|}\mathfrak b_n(C)
    =
    \sum_{\mu\in\mathcal E_{\mathbf 1}}m_A(\mu)\mu^n,
  \]
  because the average of every non-trivial irreducible character
  vanishes. Likewise, for every non-trivial \(\pi\in\Irr(G)\),
  character orthogonality gives
  \[
    \sum_{C\subset G}\frac{|C|}{|G|}\mathfrak b_n(C)\chi_\pi(C)
    =
    \sum_{\beta\in\mathcal E_\pi}m_{A_\pi}(\beta)\beta^n.
  \]
  The assumed bound \(\mathfrak b_n(C)=O(\Gamma_1^n)\) would therefore
  force each of these spectral-edge sums to be \(O(\Gamma_1^n)\).

  By definition of \(R\), at least one among the finitely many sums
  \[
    \sum_{\mu\in\mathcal E_{\mathbf 1}}m_A(\mu)\mu^n,
    \qquad
    \sum_{\beta\in\mathcal E_\pi}m_{A_\pi}(\beta)\beta^n
    \quad (\pi\neq\mathbf 1),
  \]
  is a non-zero finite exponential sum whose frequencies all have
  modulus \(R\). \Cref{lem:finite-part-exponential-sum-limsup} then
  gives limsup \(R\), contradicting the bound \(O(\Gamma_1^n)\) with
  \(\Gamma_1<R\). Hence the maximum cannot be smaller than \(R\), and
  equality follows.
\end{proof}
